\section{Related work}\label{sec:related}
\paragraph{Packing lower bounds.} van Bevern, Froese and Komusiewicz~\cite{BFK18} parameterize $F$-free editing above the size of a vertex-disjoint packing; for Cluster Editing they obtain an $O(4^\ell\ell^3+m+n)$-time algorithm for $P_3$ packings (Thm.~6.13), and for cost-$t$ packings an $O(1.62^{(2t+1)\ell}+nm+n1.62^t)$ algorithm and a kernel with $O(t\ell)$ vertices. They show NP-hardness at $\ell=0$ of $K_3$-free editing for edge-disjoint packings and of $K_q$-free editing, $q\ge6$, for vertex-disjoint packings, and leave open Cluster Editing above edge-disjoint packings and the stronger modification-disjoint bound. Li, Pilipczuk and Sorge~\cite{LPS23} answer the latter: NP-hard at $\ell=0$ with at most $23$ packed $P_3$'s per vertex, $n^{2\ell+O(1)}$ for $c=2$, and open for $c=2$ in $\ell$ and for $3\le c\le22$ at $\ell=0$. They report ``grid-like structures'' as the obstacle. Our $\tw(I(H))$ parameter does not appear in either paper.

\paragraph{Twins and normal forms.} The twin lemma (Lemma~\ref{lem:twin}) is the critical-clique property behind the kernelizations of Guo~\cite{Guo09}; the closed-neighbourhood reduction rule that generalises it is Reduction Rule~6.2 of~\cite{BFK18}, which Section~\ref{sec:K} applies to components of $I(H)$. Pair-maximality and the structure of $R$ (Lemma~\ref{lem:R}) are our own bookkeeping; we make no priority claim.

\paragraph{Reduction rule on components and unions of members.} Section~\ref{sec:reduction} applies Reduction Rule~6.2 of~\cite{BFK18} to components of $I(H)$ of any size (Theorem~\ref{thm:three}); its soundness is Lemma~6.3 of~\cite{BFK18}, which we cite and do not re-prove (it was supported by our own tests on $850$ applications). The counting argument for the number of irreducible components follows the structure of the argument of~\cite{BFK18} but is carried out here with Lemma~\ref{lem:D}. The family $N^*(F)$ of Proposition~\ref{prop:five} only shows that the rule applied to pair-closed unions of members does not reduce a family with $w=1$; it says nothing about the rule for $c=1$.

\paragraph{Treewidth of expanders.} Propositions~\ref{prop:expander} and~\ref{prop:five} use the fact that cubic expanders have treewidth linear in the number of vertices, which we cite from Grohe and Marx~\cite{GM09} (their relation between tree width, bramble size and expansion); we did not re-derive or re-check it.

\paragraph{Citations and treewidth.} Of the citing works recorded by the notes, we opened the survey of Crespelle et al.~\cite{Crespelle20}, the vertex-splitting paper~\cite{Split23} (it uses modification-disjoint $P_3$ packings as a lower-bound tool), Bocci et al.~\cite{Bocci22} and Garvardt et al.~\cite{Garvardt23}; none improves $c=2$, settles $3\le c\le22$, or uses a treewidth-type parameter. Two further citing items (the PACE 2021 report and a 2021 thesis) were not opened, and citation coverage was incomplete (Semantic Scholar listed seven records, OpenAlex a count of two); this is a ``not found'', not proof of absence. We found no paper stating the complexity of plain Cluster Editing or Deletion parameterized by treewidth. The nearby $s$-Club Cluster Edge Deletion is FPT in $s+\tw$ but W[1]-hard in pathwidth~\cite{Club25}; as a deletion problem with a diameter bound it does not transfer to editing. Komusiewicz and Uhlmann~\cite{KU12} show that Cluster Editing is NP-hard at maximum degree $6$, with no $2^{o(k)}$ algorithm under ETH even at bounded degree; this is consistent with the fact that bounded degree of $I(H)$ does not give tractability (Theorem~\ref{thm:eth}).
